% !TEX root = task-4-2-2.tex
% Editable task bank. Append new BankTask blocks at the end of the document.
\documentclass[12pt,a4paper]{article}
\usepackage{physics-textbook}
\usepackage{graphicx,needspace}
\usetikzlibrary{decorations.pathmorphing}
\newif\ifShowSolutions
\ifdefined\TutorEdition
  \ShowSolutionstrue
\else
  \ShowSolutionsfalse % Change to \ShowSolutionstrue to display solutions.
\fi
\newcounter{banktask}
\newsavebox{\BankTaskBox}
\newcommand{\BankTaskContent}[3]{%
  \noindent\textbf{\thebanktask.}\quad #2\par
  \ifShowSolutions
    \Needspace{50mm}% Keep the solution heading with its opening text or diagram.
    #3\par
    \ifdefined\TutorEdition
      \ifcsname TutorSource#1\endcsname
        \nopagebreak\smallskip{\small\csname TutorSource#1\endcsname\par}%
      \fi
    \fi
  \fi
}
\newcommand{\BankTask}[3]{%
  \refstepcounter{banktask}\par\medskip
  \sbox{\BankTaskBox}{\begin{minipage}{\linewidth}%
    \BankTaskContent{#1}{#2}{#3}\end{minipage}}%
  \ifdim\dimexpr\ht\BankTaskBox+\dp\BankTaskBox\relax>\dimexpr\textheight-20mm\relax
    \Needspace{35mm}\BankTaskContent{#1}{#2}{#3}%
  \else
    \Needspace{\dimexpr\ht\BankTaskBox+\dp\BankTaskBox\relax}%
    \noindent\usebox{\BankTaskBox}\par
  \fi
  \medskip
}
\begin{document}
\BookSubsection{4.2.2}{Плоское зеркало}


% Source sdamgia-5368: https://phys-ege.sdamgia.ru/problem?id=5368
\BankTask{sdamgia-5368}{
Предмет находится на расстоянии $60\,\text{см}$ от плоского зеркала. Каково будет расстояние между ним и его изображением, если предмет приблизить к зеркалу на $25\,\text{см}$? (в сантиметрах.)
}{
\smallskip\noindent\textbf{Решение}\par
Расстояние между предметом и его изображением в плоском зеркале равно удвоенному расстоянию до предмета и равно $( 60 - 25 ) \cdot 2=70\,\text{см}.$

Ответ: $70$.
}

% Source sdamgia-6897: https://phys-ege.sdamgia.ru/problem?id=6897
\BankTask{sdamgia-6897}{
Точечный источник света находится на расстоянии $1{,}2\,\text{м}$ от плоского зеркала. На сколько уменьшится расстояние между источником и его изображением, если, не поворачивая зеркала, пододвинуть его ближе к источнику на $0{,}3\,\text{м}$? (в метрах.)
}{
\smallskip\noindent\textbf{Решение}\par
Расстояние между предметом и его изображением в плоском зеркале равно удвоенному расстоянию до предмета.
В первом случае расстояние между источником и изображением равно: $( 1{,}2 ) \cdot 2=2{,}4\,\text{м}.$
После того как зеркало пододвинули: $( 1{,}2 - 0{,}3 ) \cdot 2=1{,}8\,\text{м}.$
Таким образом, расстояние уменьшилось на $2{,}4 - 1{,}8=0{,}6 \,\text{м}.$

Ответ: $0{,}6$.
}

% Source sdamgia-31888: https://phys-ege.sdamgia.ru/problem?id=31888
\BankTask{sdamgia-31888}{
Расстояние между небольшим предметом и плоским зеркалом равно $6\,\text{см}$. Во сколько раз увеличится расстояние между предметом и его изображением, если расстояние от предмета до зеркала увеличить на $3\,\text{см}$?
}{
\smallskip\noindent\textbf{Решение}\par
Учтем, что расстояние от предмета до зеркала равно расстоянию от зеркала до мнимого изображения. Тогда первоначально расстояние между предметом и изображением было равно $12\,\text{см}$. При увеличении расстояния от предмета до зеркала его значение будет равно $9\,\text{см}$. Тогда между предметом и изображением расстояние станет равным $18\,\text{см}$ и увеличится в $1{,}5$ раза.

Ответ: $1{,}5$.
}

% Source sdamgia-35173: https://phys-ege.sdamgia.ru/problem?id=35173
\BankTask{sdamgia-35173}{
Предмет поставлен на $0{,}9\,\text{м}$ от зеркала. На сколько метров нужно приблизить предмет к зеркалу, чтобы расстояние от предмета до изображение уменьшилось в $3$ раза?
}{
\smallskip\noindent\textbf{Решение}\par
Поскольку расстояние от предмета до зеркала равно расстоянию от зеркала до изображения, то расстояние от предмета до изображения равно $1{,}8\,\text{м}$. При уменьшении этой величины в $3$ раза данное расстояние станет равным $0{,}6\,\text{м}$. Следовательно, расстояние от предмета до зеркала станет равным $0{,}3\,\text{м}$. Тогда предмет нужно придвинуть к зеркалу на расстояние $0{,}9 - 0{,}3=0{,}6\,\text{м}.$

Ответ: $0{,}6$.
}

% Source sdamgia-1739: https://phys-ege.sdamgia.ru/problem?id=1739
\BankTask{sdamgia-1739}{
\begin{center}\begin{tikzpicture}[book diagram,>=stealth]\draw[task grid,step=.5](-2,-1.5)grid(2,2);\draw[very thick](0,0)--(0,1);\foreach\y in{0,.15,...,.9}\draw(0,\y)--(-.12,\y-.12);\node[above left]at(0,1){З};\draw[->,thick](1,0)--(1,-1)node[midway,right]{С};\fill(1,0.5)circle(.04);\node[right]at(1,0.5){$\Gamma$};\end{tikzpicture}\end{center}
В плоском зеркале $3$ наблюдается изображение стрелки С, глаз находится в точке Г. Какая часть (доля) изображения стрелки в зеркале видна глазу?
}{
\smallskip\noindent\textbf{Решение}\par
\begin{center}\begin{tikzpicture}[book diagram,>=stealth]\draw[task grid,step=.5](-2,-1.5)grid(2,2);\draw[very thick](0,0)--(0,1);\foreach\y in{0,.15,...,.9}\draw(0,\y)--(-.12,\y-.12);\node[above left]at(0,1){З};\draw[->,thick](1,0)--(1,-1)node[midway,right]{С};\fill(1,0.5)circle(.04);\node[right]at(1,0.5){$\Gamma$};\draw[dashed](1,0.5)--(0,0);\draw(0,0)--(-1.5,-0.75);\draw[dashed](1,0.5)--(0,1);\draw(0,1)--(-1.5,1.75);\draw[->,thick](-1,0)--(-1,-1);\end{tikzpicture}\end{center}

На рисунке построено изображение стрелки C в плоском зеркале и обозначена область, которая видна глазу в зеркале из точки Г. Из рисунка ясно, что глазу видна только половина стрелки.

Ответ: $0{,}5$.
}

% Source sdamgia-1740: https://phys-ege.sdamgia.ru/problem?id=1740
\BankTask{sdamgia-1740}{
\begin{center}\begin{tikzpicture}[book diagram,>=stealth]\draw[task grid,step=.5](-2,-1.5)grid(2,2);\draw[very thick](0,0)--(0,1);\foreach\y in{0,.15,...,.9}\draw(0,\y)--(-.12,\y-.12);\node[above left]at(0,1){З};\draw[->,thick](1,0)--(1,-1)node[midway,right]{С};\fill(1,0.5)circle(.04);\node[right]at(1,0.5){$\Gamma$};\end{tikzpicture}\end{center}
В плоском зеркале $3$ наблюдается изображение стрелки С, глаз находится в точке Г. На сколько клеток нужно сместить глаз по вертикали, чтобы полностью увидеть изображение стрелки? (Смещение считайте положительным при движении вверх и отрицательным при движении вниз.)
}{
\smallskip\noindent\textbf{Решение}\par
\begin{center}\begin{tikzpicture}[book diagram,>=stealth]
\draw[task grid,step=.5](-2,-2)grid(2,2);
\draw[orange!85!black,very thick](0,0)--(0,1);
\foreach \y in {0,.25,.5,.75}{\draw[orange!85!black,thick] (0,\y)--(-.15,\y+.15)--(-.15,\y);}
\node[above left]at(0,1){З};
\draw[->,thick](1,0)--(1,-1)node[midway,right]{С};
\draw[->,thick,red!82!black](-1,0)--(-1,-1);
\draw[thick,blue!68!black](-1.5,-.75)--(0,0) (-1,1.5)--(0,1);
\draw[dashed,thick,blue!68!black](0,0)--(1,.5) (0,1)--(1,.5);
\draw[thick,green!55!black](-1.5,1)--(1,1) (-1.5,-1.5)--(1,1);
\fill(1,.5)circle(.046)node[above right]{$\Gamma$};
\fill[green!55!black](1,1)circle(.046);
\end{tikzpicture}\end{center}

На рисунке построено изображение стрелки C в плоском зеркале и обозначена область, которая видна глазу в зеркале из точки Г. Перемещение глаза на одну клетку вверх позволит полностью увидеть изображение стрелки.

Ответ: $1$.
}

% Source sdamgia-1711: https://phys-ege.sdamgia.ru/problem?id=1711
\BankTask{sdamgia-1711}{
Ученики исследовали соотношение между скоростями автомобильчика и его изображения в плоском зеркале в системе отсчета, связанной с зеркалом (см. рис.). 
\begin{center}\begin{tikzpicture}[book diagram,>=stealth]\draw[->](-2.5,0)--(2.5,0)node[right]{$x$};\draw[->](0,-.2)--(0,1.5)node[right]{$y$};\draw[very thick](0,0)--(0,1.3);\node[below]at(0,0){$0$};\draw(.7,.1)--(.7,.3)--(.95,.3)--(1.1,.6)--(1.5,.6)--(1.7,.3)--(1.9,.25)--(1.9,.1)--cycle;\draw(1,.1)circle(.12) (1.6,.1)circle(.12);\draw[->](1.7,.4)--(2.3,.4)node[above]{$\vec v$};\begin{scope}[xscale=-1,dashed]\draw(.7,.1)--(.7,.3)--(.95,.3)--(1.1,.6)--(1.5,.6)--(1.7,.3)--(1.9,.25)--(1.9,.1)--cycle;\draw(1,.1)circle(.12) (1.6,.1)circle(.12);\end{scope}\end{tikzpicture}\end{center}
Проекция на ось Ох вектора скорости, с которой движется изображение, в этой системе отсчета равна

$1$)  $- 2 v$ 
$2$)  $2 v$ 
$3$)  $v$ 
$4$)  $- v$
}{
\smallskip\noindent\textbf{Решение}\par
Изображение в плоском зеркале расположено симметрично объекту относительно зеркала. Изображение удаляется от зеркала со скоростью $- \vec v .$ Проекция на ось Ox вектора скорости изображения равна $- v .$
Ответ: $4$.
}

% Source sdamgia-3590: https://phys-ege.sdamgia.ru/problem?id=3590
\BankTask{sdamgia-3590}{
Точечный источник S расположен вблизи системы, состоящей из двух плоских зеркал $\,\text{З}_1$ и $\,\text{З}_2,$ так, как показано на рисунке. Сколько изображений даст эта система зеркал?

\begin{center}\begin{tikzpicture}[book diagram,>=stealth]\draw[task grid,step=.4](0,0)grid(4.8,4.8);\draw[very thick](2.4,3.6)--(2.4,2.4)--(3.6,.8);\foreach\y in{2.4,2.6,...,3.5}\draw(2.4,\y)--++(.15,-.1);\foreach\t in{0,.2,...,1}{\draw({2.4+1.2*\t},{2.4-1.6*\t})--++(.15,0);}\fill(.8,4)circle(.04);\node[above]at(.8,4){$S$};\node[right]at(2.4,3.4){$\text{З}_1$};\node[right]at(3.6,.9){$\text{З}_2$};\end{tikzpicture}\end{center}
}{
\smallskip\noindent\textbf{Решение}\par
Изображение в плоском зеркале формирует расходящийся пучок света, отразившийся от зеркала. Глаз, продолжая лучи «за зеркало», видит там мнимое изображение. При этом для того, чтобы изображение в плоском зеркале существовало, вовсе не обязательно, чтобы предмет располагался непосредственно напротив зеркала. Для того достаточно, чтобы источник располагался в полупространстве, к которому обращена отражающая поверхность зеркала. Например, точки А и Б дают изображение в зеркале, а точка В  — нет.

\begin{center}\begin{tikzpicture}[book diagram,>=stealth]\draw[thick](-1.8,0)--(2.2,0);\foreach\x in{-1.8,-1.55,...,2.1}\draw(\x,0)--++(-.12,-.15);\fill(-1,1.8)circle(.04)(2,1.1)circle(.04)(-2,-1)circle(.04)(-1,-1.8)circle(.04)(2,-1.1)circle(.04);\node[above]at(-1,1.8){А};\node[right]at(2,1.1){Б};\node[below]at(-2,-1){В};\draw[->](-1,1.8)--(-0.2,0);\draw[->](-0.2,0)--(0.52000,1.62000);\draw[dashed](-0.2,0)--(-1,-1.8);\draw[->](-1,1.8)--(0.5,0);\draw[->](0.5,0)--(1.85000,1.62000);\draw[dashed](0.5,0)--(-1,-1.8);\draw[->](2,1.1)--(-1.3,0);\draw[->](-1.3,0)--(-4.27000,0.99000);\draw[dashed](-1.3,0)--(2,-1.1);\draw[->](2,1.1)--(0.4,0);\draw[->](0.4,0)--(-1.04000,0.99000);\draw[dashed](0.4,0)--(2,-1.1);\end{tikzpicture}\end{center}

Из рисунка видно, что источник S располагается перед обоими зеркалами, значит, он будет давать изображения в обоих зеркалах. Поскольку зеркала расположены таким образом, что лучи, отразившиеся от одного зеркала, на другое не попадают, заключаем, что изображения, получающегося после двукратного отражения, наблюдаться не может. Следовательно, данная система зеркал дает $2$ изображения.

\begin{center}\begin{tikzpicture}[book diagram,>=stealth]\draw[task grid,step=.4](0,0)grid(4.8,4.8);\draw[very thick](2.4,3.6)--(2.4,2.4)--(3.6,.8);\foreach\y in{2.4,2.6,...,3.5}\draw(2.4,\y)--++(.15,-.1);\foreach\t in{0,.2,...,1}{\draw({2.4+1.2*\t},{2.4-1.6*\t})--++(.15,0);}\fill(.8,4)circle(.04);\node[above]at(.8,4){$S$};\node[right]at(2.4,3.4){$\text{З}_1$};\node[right]at(3.6,.9){$\text{З}_2$};\fill(4,4)circle(.035);\node[above]at(4,4){$S_1$};\fill(1.312,4.384)circle(.035);\node[above]at(1.312,4.384){$S_2$};\draw[->](.8,4)--(2.4,2.8);\draw[->](2.4,2.8)--(1.44,2.0799999999999996);\draw[dashed](2.4,2.8)--(4,4);\draw[->](.8,4)--(2.4,3.2);\draw[->](2.4,3.2)--(1.44,2.72);\draw[dashed](2.4,3.2)--(4,4);\draw[->](.8,4)--(2.6,2.13333);\draw[->](2.6,2.13333)--(3.3728000000000002,0.7829279999999998);\draw[dashed](2.6,2.13333)--(1.312,4.384);\draw[->](.8,4)--(3.2,1.33333);\draw[->](3.2,1.33333)--(4.332800000000001,-0.4970720000000002);\draw[dashed](3.2,1.33333)--(1.312,4.384);\end{tikzpicture}\end{center}

Ответ: $2$.
}

% Source sdamgia-9092: https://phys-ege.sdamgia.ru/problem?id=9092
\BankTask{sdamgia-9092}{
На рисунке изображены два квадратных плоских зеркала, касающиеся друг друга краями (см. рис. слева). Угол раствора зеркал $90^\circ$. На линии OO', проходящей через линию касания зеркал перпендикулярно к ней, помещен точечный источник света S. Точки $S_{1}$, $S_{2}$ и $S_{3}$  — изображения источника в этих зеркалах при данном угле раствора. Угол раствора зеркал увеличивают до $120^\circ$ (см. рис. справа).

\begin{center}\begin{tikzpicture}[book diagram,>=stealth]\draw[dashed](-2.7,0)--(2.7,0);\node[left]at(-2.7,0){$O$};\node[right]at(2.7,0){$O^\prime$};\fill(2,0)circle(.04);\node[above]at(2,0){$S$};\draw[thick](0,0)--(1.5556349186104048,-1.5556349186104048);\draw[dashed](0,0)--(-1.5556349186104048,1.5556349186104048);\draw(0.10607,-0.10607)--(-0.06393,-0.10607);\draw(0.24749,-0.24749)--(0.07749,-0.24749);\draw(0.38891,-0.38891)--(0.21891,-0.38891);\draw(0.53033,-0.53033)--(0.36033,-0.53033);\draw(0.67175,-0.67175)--(0.50175,-0.67175);\draw(0.81317,-0.81317)--(0.64317,-0.81317);\draw(0.95459,-0.95459)--(0.78459,-0.95459);\draw(1.09602,-1.09602)--(0.92602,-1.09602);\draw(1.23744,-1.23744)--(1.06744,-1.23744);\draw(1.37886,-1.37886)--(1.20886,-1.37886);\draw(1.52028,-1.52028)--(1.35028,-1.52028);\draw[thick](0,0)--(1.5556349186104048,1.5556349186104048);\draw[dashed](0,0)--(-1.5556349186104048,-1.5556349186104048);\draw(0.10607,0.10607)--(-0.06393,0.10607);\draw(0.24749,0.24749)--(0.07749,0.24749);\draw(0.38891,0.38891)--(0.21891,0.38891);\draw(0.53033,0.53033)--(0.36033,0.53033);\draw(0.67175,0.67175)--(0.50175,0.67175);\draw(0.81317,0.81317)--(0.64317,0.81317);\draw(0.95459,0.95459)--(0.78459,0.95459);\draw(1.09602,1.09602)--(0.92602,1.09602);\draw(1.23744,1.23744)--(1.06744,1.23744);\draw(1.37886,1.37886)--(1.20886,1.37886);\draw(1.52028,1.52028)--(1.35028,1.52028);\draw[dashed](2,0)--(0.00000,-2.00000);\fill(0.00000,-2.00000)circle(.035);\node[below]at(0.00000,-2.00000){$ S_2 $};\draw[dashed](2,0)--(0.00000,2.00000);\fill(0.00000,2.00000)circle(.035);\node[above]at(0.00000,2.00000){$ S_1 $};\draw[dashed](0,2)--(-2,0)--(0,-2);\fill(-2,0)circle(.035);\node[above]at(-2,0){$S_3$};\begin{scope}[shift={(1,1)},rotate=45]\draw(.18,0)--(.18,.18)--(0,.18);\end{scope}\begin{scope}[shift={(1,-1)},rotate=-45]\draw(.18,0)--(.18,.18)--(0,.18);\end{scope}\begin{scope}[shift={(-1,1)},rotate=-45]\draw(.18,0)--(.18,.18)--(0,.18);\end{scope}\begin{scope}[shift={(-1,-1)},rotate=45]\draw(.18,0)--(.18,.18)--(0,.18);\end{scope}\end{tikzpicture}\end{center}

\begin{center}\begin{tikzpicture}[book diagram,>=stealth]\draw[dashed](-2.7,0)--(2.7,0);\node[left]at(-2.7,0){$O$};\node[right]at(2.7,0){$O^\prime$};\fill(2,0)circle(.04);\node[above]at(2,0){$S$};\draw[thick](0,0)--(1.1000000000000003,-1.905255888325765);\draw[dashed](0,0)--(-1.1000000000000003,1.905255888325765);\draw(0.07500,-0.12990)--(-0.09500,-0.12990);\draw(0.17500,-0.30311)--(0.00500,-0.30311);\draw(0.27500,-0.47631)--(0.10500,-0.47631);\draw(0.37500,-0.64952)--(0.20500,-0.64952);\draw(0.47500,-0.82272)--(0.30500,-0.82272);\draw(0.57500,-0.99593)--(0.40500,-0.99593);\draw(0.67500,-1.16913)--(0.50500,-1.16913);\draw(0.77500,-1.34234)--(0.60500,-1.34234);\draw(0.87500,-1.51554)--(0.70500,-1.51554);\draw(0.97500,-1.68875)--(0.80500,-1.68875);\draw(1.07500,-1.86195)--(0.90500,-1.86195);\draw[thick](0,0)--(1.1000000000000003,1.905255888325765);\draw[dashed](0,0)--(-1.1000000000000003,-1.905255888325765);\draw(0.07500,0.12990)--(-0.09500,0.12990);\draw(0.17500,0.30311)--(0.00500,0.30311);\draw(0.27500,0.47631)--(0.10500,0.47631);\draw(0.37500,0.64952)--(0.20500,0.64952);\draw(0.47500,0.82272)--(0.30500,0.82272);\draw(0.57500,0.99593)--(0.40500,0.99593);\draw(0.67500,1.16913)--(0.50500,1.16913);\draw(0.77500,1.34234)--(0.60500,1.34234);\draw(0.87500,1.51554)--(0.70500,1.51554);\draw(0.97500,1.68875)--(0.80500,1.68875);\draw(1.07500,1.86195)--(0.90500,1.86195);\end{tikzpicture}\end{center}

Определите, как при этом изменятся следующие величины: количество изображений источника в зеркалах; расстояние от источника до ближайшего к нему изображения.

$1$.  Увеличится.
$2$.  Уменьшится.
$3$.  Не изменится.

\TaskChangeTable{Количество изображений источника в зеркалах}{Расстояние от источника до ближайшего к нему изображения}
}{
\smallskip\noindent\textbf{Решение}\par
На рисунке показано, какое количество изображений получится, если угол станет $120^\circ$. Изображений источника станет $2$. Расстояние от источника до ближайшего к нему изображения увеличится.

\begin{center}\begin{tikzpicture}[book diagram,>=stealth]\draw[dashed](-2.7,0)--(2.7,0);\node[left]at(-2.7,0){$O$};\node[right]at(2.7,0){$O^\prime$};\fill(2,0)circle(.04);\node[above]at(2,0){$S$};\draw[thick](0,0)--(1.5556349186104048,-1.5556349186104048);\draw[dashed](0,0)--(-1.5556349186104048,1.5556349186104048);\draw(0.10607,-0.10607)--(-0.06393,-0.10607);\draw(0.24749,-0.24749)--(0.07749,-0.24749);\draw(0.38891,-0.38891)--(0.21891,-0.38891);\draw(0.53033,-0.53033)--(0.36033,-0.53033);\draw(0.67175,-0.67175)--(0.50175,-0.67175);\draw(0.81317,-0.81317)--(0.64317,-0.81317);\draw(0.95459,-0.95459)--(0.78459,-0.95459);\draw(1.09602,-1.09602)--(0.92602,-1.09602);\draw(1.23744,-1.23744)--(1.06744,-1.23744);\draw(1.37886,-1.37886)--(1.20886,-1.37886);\draw(1.52028,-1.52028)--(1.35028,-1.52028);\draw[thick](0,0)--(1.5556349186104048,1.5556349186104048);\draw[dashed](0,0)--(-1.5556349186104048,-1.5556349186104048);\draw(0.10607,0.10607)--(-0.06393,0.10607);\draw(0.24749,0.24749)--(0.07749,0.24749);\draw(0.38891,0.38891)--(0.21891,0.38891);\draw(0.53033,0.53033)--(0.36033,0.53033);\draw(0.67175,0.67175)--(0.50175,0.67175);\draw(0.81317,0.81317)--(0.64317,0.81317);\draw(0.95459,0.95459)--(0.78459,0.95459);\draw(1.09602,1.09602)--(0.92602,1.09602);\draw(1.23744,1.23744)--(1.06744,1.23744);\draw(1.37886,1.37886)--(1.20886,1.37886);\draw(1.52028,1.52028)--(1.35028,1.52028);\draw[dashed](2,0)--(0.00000,-2.00000);\fill(0.00000,-2.00000)circle(.035);\node[below]at(0.00000,-2.00000){$ S_2 $};\draw[dashed](2,0)--(0.00000,2.00000);\fill(0.00000,2.00000)circle(.035);\node[above]at(0.00000,2.00000){$ S_1 $};\draw[dashed](0,2)--(-2,0)--(0,-2);\fill(-2,0)circle(.035);\node[above]at(-2,0){$S_3$};\begin{scope}[shift={(1,1)},rotate=45]\draw(.18,0)--(.18,.18)--(0,.18);\end{scope}\begin{scope}[shift={(1,-1)},rotate=-45]\draw(.18,0)--(.18,.18)--(0,.18);\end{scope}\begin{scope}[shift={(-1,1)},rotate=-45]\draw(.18,0)--(.18,.18)--(0,.18);\end{scope}\begin{scope}[shift={(-1,-1)},rotate=45]\draw(.18,0)--(.18,.18)--(0,.18);\end{scope}\end{tikzpicture}\end{center}

\begin{center}\begin{tikzpicture}[book diagram,>=stealth]\draw[dashed](-2.7,0)--(2.7,0);\node[left]at(-2.7,0){$O$};\node[right]at(2.7,0){$O^\prime$};\fill(2,0)circle(.04);\node[above]at(2,0){$S$};\draw[thick](0,0)--(1.1000000000000003,-1.905255888325765);\draw[dashed](0,0)--(-1.1000000000000003,1.905255888325765);\draw(0.07500,-0.12990)--(-0.09500,-0.12990);\draw(0.17500,-0.30311)--(0.00500,-0.30311);\draw(0.27500,-0.47631)--(0.10500,-0.47631);\draw(0.37500,-0.64952)--(0.20500,-0.64952);\draw(0.47500,-0.82272)--(0.30500,-0.82272);\draw(0.57500,-0.99593)--(0.40500,-0.99593);\draw(0.67500,-1.16913)--(0.50500,-1.16913);\draw(0.77500,-1.34234)--(0.60500,-1.34234);\draw(0.87500,-1.51554)--(0.70500,-1.51554);\draw(0.97500,-1.68875)--(0.80500,-1.68875);\draw(1.07500,-1.86195)--(0.90500,-1.86195);\draw[thick](0,0)--(1.1000000000000003,1.905255888325765);\draw[dashed](0,0)--(-1.1000000000000003,-1.905255888325765);\draw(0.07500,0.12990)--(-0.09500,0.12990);\draw(0.17500,0.30311)--(0.00500,0.30311);\draw(0.27500,0.47631)--(0.10500,0.47631);\draw(0.37500,0.64952)--(0.20500,0.64952);\draw(0.47500,0.82272)--(0.30500,0.82272);\draw(0.57500,0.99593)--(0.40500,0.99593);\draw(0.67500,1.16913)--(0.50500,1.16913);\draw(0.77500,1.34234)--(0.60500,1.34234);\draw(0.87500,1.51554)--(0.70500,1.51554);\draw(0.97500,1.68875)--(0.80500,1.68875);\draw(1.07500,1.86195)--(0.90500,1.86195);\draw[dashed](2,0)--(-1.00000,-1.73205);\fill(-1.00000,-1.73205)circle(.035);\node[below]at(-1.00000,-1.73205){$ S^{\prime\prime} $};\draw[dashed](2,0)--(-1.00000,1.73205);\fill(-1.00000,1.73205)circle(.035);\node[above]at(-1.00000,1.73205){$ S^\prime $};\draw(0:.55)arc(0:60:.55);\node at(30:.9){$60^\circ$};\draw(120:.55)arc(120:180:.55);\node at(150:.9){$60^\circ$};\begin{scope}[shift={(.5,.866025)},rotate=60]\draw(.18,0)--(.18,.18)--(0,.18);\end{scope}\end{tikzpicture}\end{center}

Ответ: $21$.
}

% Source sdamgia-3901: https://phys-ege.sdamgia.ru/problem?id=3901
\BankTask{sdamgia-3901}{
\begin{center}\begin{tikzpicture}[book diagram,>=stealth]\draw[thick](-2,0)--(0,2)--(2,0);\foreach\t in{0,.15,...,1.95}{\draw(-2+\t,\t)--++(0,.18);\draw(2-\t,\t)--++(0,.18);}\draw[dashed](0,-.3)--(0,2.5) (-2,0)--(2,0);\draw[<->,thick](-.7,0)--(.7,0);\node[below]at(.7,0){Л};\draw[->,thick](-1.6,0)--(-1,0);\node[below]at(-1.3,0){П};\fill(0,2)circle(.04);\node[above]at(0,2){$F$};\node at(.35,1.4){$45^\circ$};\end{tikzpicture}\end{center}
Два плоских зеркала образуют прямой двугранный угол, перпендикулярно биссектрисе которого расположена небольшая собирающая линза Л, а ее фокус F находится в вершине угла (см.
рисунок). В плоскости линзы рядом с ней находится небольшой предмет П. Постройте изображение предмета, которое получится в результате двух отражений от зеркал и последующего преломления света линзой. На каком расстоянии от предмета будет находиться его изображение?
}{
\smallskip\noindent\textbf{Решение}\par
\begin{center}\begin{tikzpicture}[book diagram,>=stealth]\draw[thick](-2,0)--(0,2)--(2,0);\foreach\t in{0,.15,...,1.95}{\draw(-2+\t,\t)--++(0,.18);\draw(2-\t,\t)--++(0,.18);}\draw[dashed](0,-.3)--(0,2.5) (-2,0)--(2,0);\draw[<->,thick](-.7,0)--(.7,0);\node[below]at(.7,0){Л};\draw[->,thick](-1.6,0)--(-1,0);\node[below]at(-1.3,0){П};\fill(0,2)circle(.04);\node[above]at(0,2){$F$};\node at(.35,1.4){$45^\circ$};\draw[dashed](0,-4)--(0,4.5);\node[left]at(0,4){$2F$};\node[right]at(0,-2){$F$};\node[right]at(0,-4){$2F$};\draw[->,thick](1.6,4)--(1,4);\node[above]at(1.3,4){$\text{П}^\prime$};\draw[->,thick](-1.6,-4)--(-1,-4);\node[below]at(-1.3,-4){$\text{П}^{\prime\prime}$};\draw[dashed](-1.3,0)--(1.3,4);\draw[dashed](1.3,4)--(0,0);\draw[->](0,0)--(-1.3,-4);\draw[dashed](1.3,4)--(.5,0);\draw[->](.5,0)--(-1.3,-4);\draw[dashed](-1.3,0)--(-1.3,-4) (0,4)--(1.6,4) (-1.6,-4)--(0,-4);\draw[->,thick](-1.3,0)--(-.8,1.2);\draw[->,thick](-.8,1.2)--(.3294118,1.6705882);\draw[->,thick](.3294118,1.6705882)--(-.3666667,0);\draw[->,thick](-.3666667,0)--(-1.3,-4);\end{tikzpicture}\end{center}
Построим вначале изображение предмета П после двух отражений от плоских зеркал. Как видно из построения на рисунке, предмет как бы «отражается» относительно вершины прямого угла, и его мнимое перевернутое изображение П' находится на расстоянии 2F от плоскости линзы Л. Далее построим изображение П« «предмета» П' в линзе - оно действительное, перевернутое относительно П' (и прямое относительно предмета П) и находится также на расстоянии 2F от линзы. Таким образом, изображение П» прямое, действительное и находится на расстоянии 2F от предмета П - ниже его, на линии, параллельной биссектрисе угла из зеркал. 
Ответ: см. рис.; изображение П" прямое, действительное и находится на расстоянии 2F от предмета П - ниже его, на линии, параллельной биссектрисе угла из зеркал.
}


% Source 758809: https://ege.fipi.ru/bank/questions.php?search=1&proj=BA1F39653304A5B041B656915DC36B38&qid=758809
\BankTask{758809}{
На шахматной доске на расстоянии двух клеток от вертикального плоского зеркала стоит слон. На сколько клеток увеличится расстояние между слоном и его изображением, если фигуру на три клетки отодвинуть от зеркала?
}{
\smallskip\noindent\textbf{Решение от GPT}\par
Изображение находится на таком же расстоянии за зеркалом, как предмет перед ним, поэтому расстояние между ними равно \(D=2d\). Увеличение расстояния составит \(\Delta D=2\Delta d=2\cdot3=6\) клеток. Ответ: \(6\) клеток.
}

% Source F3157e: https://ege.fipi.ru/bank/questions.php?search=1&proj=BA1F39653304A5B041B656915DC36B38&qid=F3157e
\BankTask{F3157e}{
Точечный источник света находится на расстоянии \(75\ \text{см}\) от плоского зеркала. На сколько уменьшится расстояние между источником и его изображением, если, не поворачивая зеркала, придвинуть его к источнику на \(25\ \text{см}\) (в $\text{см}$)?
}{
\smallskip\noindent\textbf{Решение от GPT}\par
Расстояние между источником и изображением равно \(D=2d\). Первоначально \(D_1=2\cdot75=150\ \text{см}\), после перемещения зеркала \(D_2=2(75-25)=100\ \text{см}\). Уменьшение: \(D_1-D_2=50\ \text{см}\). Ответ: \(50\ \text{см}\).
}

% Source 0439A7: https://ege.fipi.ru/bank/questions.php?search=1&proj=BA1F39653304A5B041B656915DC36B38&qid=0439A7
\BankTask{0439A7}{
Точечный источник света расположен перед плоским зеркалом на расстоянии \(1{,}2\ \text{м}\) от него. На сколько необходимо приблизить зеркало, не поворачивая его, к источнику, чтобы расстояние между источником и его изображением в зеркале уменьшилось в \(2\) раза (в $\text{м}$)?
}{
\smallskip\noindent\textbf{Решение от GPT}\par
Расстояние между источником и изображением равно \(D=2d\). Чтобы оно уменьшилось в \(2\) раза, расстояние до зеркала должно стать \(d_2=d_1/2=0{,}6\ \text{м}\). Перемещение зеркала: \(x=d_1-d_2=1{,}2-0{,}6=0{,}6\ \text{м}\). Ответ: \(0{,}6\ \text{м}\).
}

\end{document}
